%% en/sections/appendici/app_storica.tex — Extended Historical Notes EN

\chapter{Extended Historical Notes}
\markboth{Extended Historical Notes}{Extended Historical Notes}

\section{Key Chronology}

\begin{tabular}{lp{10cm}}
\toprule
\cinzel Year & \cinzel Event \\
\midrule
1194  & Frederick II born at Jesi. Norman Sicily enters the Hohenstaufen sphere. \\
1250  & Death of Frederick II. End of the \textit{Stupor Mundi}. \\
1266  & Charles of Anjou defeats Manfred at Benevento. \\
1268  & Execution of Conradin in Naples. End of the Hohenstaufen dynasty. \\
1267  & Treaty of Viterbo: Charles obtains rights over Constantinople. \\
1277  & Charles purchases the title of King of Jerusalem. \\
1281  & Election of Martin IV, pro-Angevin pope. \\
1281  & Charles plans Byzantine expedition: 300 ships, spring 1282. \\
1282 Jan & Requisitions for the expedition reach unsustainable levels. \\
1282 Feb & Giovanni da Procida is (perhaps) in Palermo. The network activates. \\
1282 Mar 30 & Sicilian Vespers: uprising in Palermo, church of Santo Spirito. \\
1282 Mar 31 & The uprising spreads. Nearly all French in Palermo are killed. \\
1282 Aug & Peter III of Aragon lands in Sicily. \\
1282--1302 & War of the Vespers. \\
1285  & Death of Charles of Anjou. \\
1302  & Peace of Caltabellotta. Sicily remains Aragonese. \\
\bottomrule
\end{tabular}

\secsep

\section{Primary Sources}

\subsection{Bartholomew of Neocastro}
\textit{Historia Sicula} (c.~1293). Written by a Sicilian contemporary,
the closest source to the events. Presents the uprising as just resistance
against a foreign tyrant. Must be read with awareness of its perspective.

\subsection{Giovanni Villani}
\textit{Nuova Cronica} (14th century). Florentine, writing decades later.
Emphasises Giovanni da Procida's role as architect of the conspiracy.
His version is the most dramatic and most influential on later collective
memory.

\subsection{Saba Malaspina}
\textit{Rerum Sicularum Historia}. Pro-Angevin perspective. Valuable for
understanding how Charles of Anjou and his supporters interpreted the uprising.

\subsection{The Genoese \textit{Liber Iurium}}
Commercial and diplomatic documents showing how the uprising was perceived
in the Mediterranean's commercial centres.

\secsep

\section{Modern Historiography}

\begin{description}[leftmargin=2em, style=nextline]
  \item[Steven Runciman, \textit{The Sicilian Vespers} (1958)]
    The most complete and readable reconstruction. Broadly accepts the
    Aragonese conspiracy narrative. Still the obligatory starting point.
  \item[David Abulafia, \textit{Frederick II} (1988)]
    Not specifically about the Vespers, but essential for understanding
    the Hohenstaufen Sicily from which the uprising emerges.
  \item[Jean Dunbabin, \textit{Charles I of Anjou} (1998)]
    The best English-language study of Charles of Anjou. Scales back
    Charles's ``villainy'' in favour of a structural reading.
  \item[Henri Bresc, \textit{Un monde méditerranéen} (1986)]
    Fundamental study of medieval Sicily as an economic and cultural space.
\end{description}

\filettodoppio
